\documentclass[twoside]{article}
\usepackage{amsmath,amssymb,graphicx}
\pagestyle{empty}

%%------------------------------------------------------------
\begin{document}

%%%%%%%%% TITLE of abstract %%%%%%%%%

\begin{center}
\textbf{\large The polluted atmosphere as a shallow domain}
\end{center}

\medskip

%%%%%%%%%  SPEAKER %%%%%%%%%

\begin{center}
Nora Juhasz\\
\emph{University  of L'Aquila}\\
\texttt{ms.juhasz.nora@gmail.com}
\end{center}

\bigskip

%%%%%%%%%  text of ABSTRACT to be included below  %%%%%%
\noindent
Considered as a geophysical fluid, the polluted atmosphere shares the shallow domain characteristics with other natural large-scale fluids such as seas and oceans. This means that its domain is excessively greater horizontally than in the vertical dimension, leading to the classic hydrostatic approximation of the Navier-Stokes equations. We consider the so-called anisotropic model as a starting point, i.e. a set of equations that describe the atmosphere including the effects of pollution, where the fluid velocity function governed by the classical, three-dimensional Navier-Stokes equations is combined with a concentration function representing the pollution.  We provide a convergence theorem for the weak solutions of the inland anisotropic model towards the hydrostatic system. Our results are valid on local domains where the use of Cartesian coordinates are legitimate and effects such as the curvature of the Earth are negligible. The shallow domain concept of the atmosphere is grasped by the small variable $\epsilon$ which stands for the aspect ratio, i.e. the ratio of the characteristic depth and characteristic width. We explicitly involve this aspect ratio in the coordinates and variables that describe the phenomena, then perform a rescaling process that makes the domain independent of $\epsilon$. Not only do we rescale the velocity and diffusion terms according to the "almost two-dimensional" concept we use, but we also take advantage of an additional scaling idea for the concentration equation which is suggested by the fact that in the downwind direction it is natural to assume that the diffusion is negligible compared to downwind transport. The main two models between which we describe the convergence result are the original anisotropic model and the hydrostatic limit model. In the first one we use the Navier-Stokes equations in all three space dimensions to describe the fluid velocity, while in the latter model we arrive to a simpler equation in the vertical dimension, namely we see the hydrostatic approximation incorporated into the model. In both models we use the delta distribution approach to mathematically represent the source term, but while in the limit model we use the delta distribution itself, in the anisotropic model it is more natural to apply nascent delta functions, i.e. Gaussian impulses that approximate the limit delta distribution. The main steps of the proof for the convergence result between the two models are \textit{a)} using the classical method of extracting a priori bounds from the energy inequality, \textit{b)} passing to the limit in the linear terms, and, finally, \textit{c)} finding and using compactness results that allow us to pass to the limit in the nonlinear terms that require stronger regularities. Our additional results include an existence theorem and a generalization of the original inland model to coastal conditions.


%%%%%%%%%  end text of ABSTRACT   %%%%%%%%%%

\bigskip


%%%%%% CO-AUTHORS %%%%%%%
%\hrule
\medskip

\noindent {{\bf Joint work with:} D. Donatelli, University of L'Aquila}

\end{document}
